\documentclass[10.5pt,a5paper,twoside]{memoir}
\usepackage[german]{babel}
\usepackage[utf8]{inputenc} 

\usepackage{csquotes}

\renewcommand{\thesection}{\arabic{section}}
\setcounter{secnumdepth}{3} % Subsection mit Zähler (1.1) versehen

%% pt sans
\usepackage[T1]{fontenc}
\usepackage{PTSerif}

\usepackage{textcomp}

%%% Other Settings %%%
\sloppy
\raggedbottom

\usepackage{rotating}
%%% disable figure numbering

%%% import
\usepackage{import}

%%% captions
\usepackage{caption}

%%% graphics
\usepackage{graphicx}
\setkeys{Gin}{width=.8\textwidth} %default pics size
\usepackage{longtable}

%%%%url
\usepackage[hyphens]{url}
\urlstyle{same}
\usepackage[colorlinks, linkcolor=black,citecolor=black, urlcolor=blue,
breaklinks= true]{hyperref}
%%% remove chapter

\renewcommand*{\chaptitlefont}{\LARGE}

%%%%%%%bod test
\linespread{1.15}

%%% microtype
\usepackage{microtype}

%%% footnote margib
\footmarkstyle{\textsuperscript{#1} }    % indent the footmark 

% Hurenkinder und Schusterjungen verhindern
\clubpenalty10000
\widowpenalty10000
\displaywidowpenalty=10000

%%
% Kolumnentitel
%%
\renewcommand{\chaptermark}[1]{ \markboth{#1}{}  } % Stil der Kopfzeile zurücksetzen
\renewcommand{\sectionmark}[1]{ \markright{#1}{} } % Stil der Kopfzeile 

\renewcommand*{\chaptitlefont}{\normalfont\LARGE\flushright}

% Page style
\makepagestyle{mystyle}
\makeevenhead{mystyle}{\small\thepage}{\small \#l10j}{}
\makeoddhead{mystyle}{}{\small\slshape\leftmark}{\small\thepage}
\makeoddfoot{plain}{}{\small\thepage}{}
\makeevenhead{headings}{\small\thepage}{}{\slshape\leftmark}
\pagestyle{mystyle}

%%
% Inhaltsverzeichnis
%%
\renewcommand*{\cftchapterdotsep}{\cftdotsep}
\settocdepth{chapter}
\renewcommand{\cftchapterfont}{\normalfont}
\renewcommand{\cftchapterpagefont}{\normalfont}

\makeatletter
\DeclareRobustCommand\authortoctext[1]{%
{\addvspace{20pt}\nopagebreak\leftskip0em\relax
\rightskip \@tocrmarg\relax
\noindent\itshape#1\par\addvspace{-7pt}}}
\makeatother
\newcommand\authortoc[1]{%
  \gdef\chapterauthor{#1}%
  \addtocontents{toc}{\authortoctext{#1}}}

\makeatletter
\renewcommand{\cfttocbeforelisthook}{\pagestyle{empty}\let\ps@plain\ps@empty}
\renewcommand{\cfttocafterlisthook}{\cleardoublepage\pagestyle{headings}}
\makeatother

\newlength\drop
\makeatletter
\newcommand*\titleM{\begingroup% Misericords, T&H p 153
\setlength\drop{0.08\textheight}
\centering
\vspace*{\drop}
{\Huge\bfseries Die Bibliothek als Idee}\\[\baselineskip]
{\scshape Beiträge des Symposiums \\ -- 10 Jahre LIBREAS. Library Ideas -- \\ am ICI Kulturlabor Berlin}\\[\baselineskip]
\vfill
{\large\scshape {\small herausgegeben vom}\\{\large LIBREAS. Verein zur Förderung der bibliotheks- und informationswissenschaftlichen Kommunikation e.V.}}\par
\vfill
{\scshape Berlin, Bielefeld, Chur, Dresden, Exeter, München}\par
\vfill
{\scshape 2016}\par
%\vspace*{2\drop}
\endgroup}
%\makeatother

%%
% Dokumentenbeginn
%%
\begin{document}

\frontmatter


\pagestyle{empty}
\begin{center}
{\LARGE DIE BIBLIOTHEK ALS IDEE} 
\end{center}

\cleartorecto

\pagestyle{empty}
\titleM
\cleartoverso


\begingroup
\footnotesize
\parindent 0pt
\parskip \baselineskip
\vspace*{\fill}


Mit freundlicher Unterstützung durch

\includegraphics[width=0.5\textwidth]{astec_Logo_300dpi.jpg}

\vspace{10mm}

LIBREAS. Verein (Hrsg.) Die Bibliothek als Idee : Beiträge des Symposiums \enquote{10
Jahre LIBREAS. Library Ideas} am ICI Kulturlabor Berlin. Berlin, 2016.

\textcopyright{}
2015 / 2016 Die Autorinnen und Autoren. 

Dieses Werk ist lizenziert unter einer Creative Commons Namensnennung 3.0 Lizenz.

Satz und Layout: Redaktion LIBREAS unter Verwendung von pandoc und \LaTeX

Druck: epubli

Webausgabe: \url{http://libreas.eu/ausgabe28/}

Quellcoderepository: \url{https://github.com/libreas/ausgabe28/}

LIBREAS. Library Ideas sind: Ben Kaden, Maxi Kindling, Manuela Schulz
(Herausgeberinnen und Herausgeber), Linda Freyberg, Najko Jahn, Karsten
Schuldt, Matti Stöhr, Doreen Thiede (Redaktion).

    \endgroup
\newpage




\noindent Wir bedanken uns bei der Firma |a|S|tec| angewandte
Systemtechnik GmbH für die finanzielle Unterstützung des Drucks.

\vspace{7mm}

\noindent Unser Dank gilt überdies der Firma Citavi und dem ICI Kulturlabor Berlin,
die das Symposium unterstützt haben.

\vspace{7mm}

\noindent Ein ebenso herzlicher Dank gebührt unseren Referentinnen und Referenten und
zahlreichen Gästen, die das Symposium zu einem lebhaften Austauschort machten.
Ein besonderer Dank gilt Corinna Haas, die uns mit dem ICI Kulturlabor eine
großartige Location vermittelt und für den Erfolg des Symposiums maßgeblich
verantwortlich war. Allen Unterstützerinnen und Unterstützern, die uns vor,
während und nach dem Symposium tatkräftig halfen, danken wir ebenso herzlich:
ohne Euch hätten wir es nicht geschafft.

\vspace{7mm}

\noindent Wir widmen den Band den Mitgliedern des LIBREAS. Vereins.


\cleartorecto

\tableofcontents*

\mainmatter
\pagestyle{mystyle}

%% Editorial

\authortoc{Redaktion LIBREAS}
\chapter*[Editorial]{Editorial}
\addcontentsline{toc}{chapter}{Editorial}
\refstepcounter{chapter}
\begin{flushright}
{\large Redaktion LIBREAS}
\end{flushright}

\vspace{5mm}

\import{../00editorial/}{master.tex}

%% Karin Aleksander
\authortoc{Karin Aleksander}
\chapter*[Ist eine transdisziplinäre Bibliothek möglich?]{Ist eine transdisziplinäre Bibliothek möglich? Oder: Wie die Geschlechterforschung Idee und Ideal der Bibliothek herausfordert}
\addcontentsline{toc}{chapter}{Ist eine transdisziplinäre Bibliothek möglich? Oder: Wie die Geschlechterforschung Idee und Ideal der Bibliothek herausfordert}
\refstepcounter{chapter}
\begin{flushright}
{\large Karin Aleksander}
\end{flushright}

\vspace{5mm}

\import{../01aleksander/}{master.tex}

%% Ute Engelkenmeier
\authortoc{Ute Engelkenmeier}
\chapter*[Bibliothek und Bibliothekare in Comedy und Komödie]{Das Bild der Bibliothek und Bibliothekare in den Fernsehgenres Comedy und Komödie}
\addcontentsline{toc}{chapter}{Das Bild der Bibliothek und Bibliothekare in den Fernsehgenres Comedy und Komödie}
\refstepcounter{chapter}
\begin{flushright}
{\large Ute Engelkenmeier}
\end{flushright}

\vspace{5mm}

{\large \noindent{ABSTRACT}}

{\small\noindent Mittels der empirischen Methode der standardisierten, quantitativen
Inhaltsanalyse werden 51 Fernsehsendungen aus den Subgenres Comedy und
Komödie untersucht, in welchen Bibliotheken als Schauplatz und
Bibliothekare als Figur vorkommen. Der betrachtete Zeitraum der
Ausstrahlungen ist von Januar 2010 bis Juli 2015. Ziel der Untersuchung
ist die Klärung der Frage, ob Bibliotheken und Bibliothekare stereotyp
dargestellt werden und ob sie innerhalb der gewählten Subgenres Objekt
des Witzes sind. Untersucht werden unter anderem die Dimensionen
Erscheinungsbild und Nutzungsmotive. Die Untersuchung kommt zu dem
Schluss, dass das Erscheinungsbild von Bibliotheken überwiegend auf das
konstitutive Element der Bücherregale reduziert ist. In der Nutzung
überwiegen Motive wie Ausleihen, Lesen, aber auch Nutzung der Bibliothek
als sozialer Treffpunkt (Kennenlernen, Date) oder die Suche nach einem
ruhigen Rückzugsort. Das Bild der Bibliothekare ist im äußeren
Erscheinungsbild zwar konservativer als andere Figuren, jedoch nicht
negativ, auch sind die Rollen nicht Objekt des Witzes. Eine Tendenz zu
einer weniger konservativen Darstellung in jüngeren Produktionen ist
festzustellen.}

\newpage 

\import{../02engelkenmeier/}{master.tex}

%%% Hartmann
\authortoc{\\Frank Hartmann}
\chapter*[Dokumentation als Gegenidee zur Bibliothek]{Dokumentation als Gegenidee zur Bibliothek}
\addcontentsline{toc}{chapter}{Dokumentation als Gegenidee zur Bibliothek}
\refstepcounter{chapter}
\begin{flushright}
{\large Frank Hartmann}
\end{flushright}

\vspace{5mm}

{\large \noindent{ABSTRACT}}

{\small\noindent
	Schon zur Wende ins 20. Jahrhundert zweifelte der belgische Privatgelehrte Paul Otlet an der Zukunft des Buches und der Bibliothek. Statt dessen begann er damit, eine Dokumentation und Neuorganisation des Weltwissens anzulegen, und mittels eines Karteikartensystems (Répertoire Bibliographique Universel) zu vernetzen. Dieses Projekt eines flexiblen, abfrageorientierten Wissensbestandes in einem \enquote{Hypermedium} (Otlet) besetzte jene technologische Leerstelle, die inzwischen eine die bibliothekarische Epoche aufsprengende neue Wissenskultur der digitalen Medialität produziert hat.}

\vspace{5mm}

\import{../04hartmann/}{master.tex}

%%% wagner
\authortoc{Kirsten Wagner}
\chapter*[Die architektonische Idee der modernen Bibliothek]{Die architektonische Idee der modernen Bibliothek}
\addcontentsline{toc}{chapter}{Die architektonische Idee der modernen Bibliothek}
\refstepcounter{chapter}
\begin{flushright}
{\large Kirsten Wagner}
\end{flushright}

\vspace{5mm}

{\large \noindent{ABSTRACT}}

{\small\noindent
	
Die moderne Bibliothek geht auf einen Übergang der Saalbibliothek zur
Magazinbibliothek zurück, der sich im 19. Jahrhundert vollzog. Bevor erste
Magazinbibliotheken ab den 1850er Jahren gebaut wurden, existierte die moderne
Bibliothek in der ersten Hälfte des 19. Jahrhunderts vor allem als Idee. Es
entstand eine kaum übersehbare Zahl an Bibliotheksentwürfen und -schriften, in
denen nach einer architektonischen Form gesucht wurde, die der Expansion der
Gutenberggalaxis, der Erweiterung des Nutzerkreises und der räumlichen
Ausdifferenzierung der Bibliotheksfunktionen entsprach. Ausgehend von der
öffentlichen Bibliotheksdiskussion, wie sie auch und besonders in der
Architekturpresse geführt wurde, untersucht der Beitrag am Beispiel der
französischen Bibliothèque royale den Entwurf der modernen Bibliothek als
Zentral- bzw. Rundbau mit einem Lesesaal in seiner Mitte. In diesen Entwurf
schreiben sich in der Zeitspanne von 1790 bis 1850 vollkommen verschiedene
Wissens- und Körperordnungen ein. Ihre Verbindung besteht dort, wo sich das
absolute räumliche Zentrum der Bibliothek durch ein allsehendes Auge besetzt
zeigt.}

\vspace{5mm}

\import{../05wagner/}{master.tex}

%%%% hobohm

\authortoc{Hans-Christoph Hobohm}
\chapter*[Vom Ort zum Akteur]{Vom Ort zum Akteur. Heterotopologie + Akteur-Network-Theorie auf die Bibliothek bezogen}
\addcontentsline{toc}{chapter}{Vom Ort zum Akteur. Heterotopologie + Akteur-Network-Theorie auf die Bibliothek bezogen}
\refstepcounter{chapter}
\begin{flushright}
{\large Hans-Christoph Hobohm}
\end{flushright}

\vspace{5mm}

\import{../06hobohm/}{master.tex}

%%%% Verzeichnis der Autorinnen und Autoren

\chapter*[Verzeichnis der Autorinnen und Autoren]{Verzeichnis der Autorinnen und Autoren}
\addcontentsline{toc}{chapter}{Verzeichnis der Autorinnen und Autoren}
\refstepcounter{chapter}

\noindent \textbf{Karin Aleksander}, Dr. phil., ist Leiterin der Bibliothek des Zentrum für transdisziplinäre Geschlechterstudien, Humboldt-Universität zu Berlin.

\vspace{5mm}

\noindent \textbf{Ute Engelkenmeier}, ausgestattet mit einem ausgiebigen Studium des
Bibliothekswesens (FHBD Köln, Musikbibliothekarisches Zusatzstudium HBI
Stuttgart) und der Bibliotheks- und Informationswissenschaft (Humboldt-Universität zu Berlin) seit 1995 in der Universitätsbibliothek Dortmund
tätig, dort zehn Jahre lang Öffentlichkeitsarbeiterin und seit vier
Jahren Geschäftsbereichsleitung Medienbereitstellung und Information
(Benutzung).
Interessensschwerpunkte: Nutzerforschung, Lernortentwicklung, Marketing
und Management.

\vspace{5mm}

\noindent \textbf{Frank Hartmann} ist Professor für Geschichte und Theorie der Visuellen Kommunikation an der Bauhaus-Universität Weimar. Er hat zahlreiche Publikationen zu Medienphilosophie verfasst und ist Herausgeber der Schriftenreihe \enquote{Forschung Visuelle Kultur}, darunter der Band \enquote{Vom Buch zur Datenbank. Paul Otlets Utopie der Wissensvisualisierung} (Berlin 2012).

\vspace{5mm}

\noindent \textbf{Hans-Christoph Hobohm} ist seit 1995 Professor für Biblio\-theks- und Informationswissenschaft an der Fachhochschule Potsdam. Neben Wissensmanagement und Informationsverhaltensforschung ist eines seiner Lehr- und Forschungsgebiete Bibliothekstypologie und Bibliothekstheorie.

\vspace{5mm}

\noindent \textbf{Kirsten Wagner} ist Professorin für Kulturwissenschaft und
Kommunikationswissenschaft am Fachbereich Gestaltung der Fachhochschule
Bielefeld. Davor war sie wissenschaftliche Mitarbeiterin am
Kulturwissenschaftlichen Institut der Humboldt-Universität zu Berlin, wo
sie im Rahmen des Sonderforschungsbereiches 447 „Kulturen des
Perfomativen“ zu historischen und gegenwärtigen Formen räumlicher
Wissensorganisation gearbeitet hat, darunter auch zur Bibliothek als
materiellem Raum des Wissens. Ihre aktuellen Arbeits- und
Forschungsschwerpunkte umfassen neben Fragen der Wissensorganisation
Raum- und Wahrnehmungstheorien sowie eine Anthropologie der Architektur.

\cleartoverso

\cleartorecto


\end{document}